\documentclass[11pt]{article}

\usepackage[margin=1in]{geometry}
\usepackage{amsmath,amssymb,amsfonts}
\usepackage{booktabs}
\usepackage{enumitem}
\usepackage{hyperref}
\usepackage{xcolor}
\usepackage{hyperref}

\hypersetup{
    colorlinks=true,
    linkcolor=blue,
    citecolor=blue,
    urlcolor=blue
}

\setlist[itemize]{leftmargin=1.5em, itemsep=0.2em, topsep=0.2em}
\setlist[enumerate]{leftmargin=1.7em, itemsep=0.25em, topsep=0.25em}

\title{Offline Metrics for Early Stopping in VLA Fine-Tuning}
\author{}
\date{}

\begin{document}
\maketitle

\section{Goal}

We want to choose an early-stopping checkpoint for fine-tuning a Vision--Language--Action (VLA) policy without running expensive real-world robot evaluation at every checkpoint.

For each checkpoint, we compute offline validation metrics and ask:
\begin{equation}
\textit{Which metric best predicts the checkpoint with the best downstream execution quality?}
\end{equation}

The candidate metrics are:
\begin{enumerate}
    \item Dynamic Time Warping distance, denoted $D_{\mathrm{DTW}}$;
    \item Optimal Transport or Sinkhorn distance, denoted $D_{\mathrm{OT}}$;
    \item Cosine distance, denoted $D_{\mathrm{cos}}$;
    \item negative log-likelihood or a model-specific likelihood surrogate, denoted $\mathrm{NLL}$.
\end{enumerate}

There are two evaluation settings:
\begin{enumerate}
    \item \textbf{Action-chunk evaluation}: compare predicted action chunks with held-out demonstration action chunks.
    \item \textbf{State-action rollout evaluation}: use a world model to roll out the policy, then compare generated state-action chunks with held-out successful state-action chunks.
\end{enumerate}

\section{Notation}

Let
\begin{equation}
\mathcal{D}=\{\tau_i\}_{i=1}^{N}
\end{equation}
be a dataset of robot demonstration trajectories. A trajectory is
\begin{equation}
\tau_i = (s_{i,0},h_{i,0},a_{i,0},s_{i,1},h_{i,1},a_{i,1},\ldots,s_{i,T_i}).
\end{equation}

The symbols are defined as follows.

\begin{center}
\begin{tabular}{ll}
\toprule
Symbol & Meaning \\
\midrule
$s_t$ & environment or robot state \\
$o_t$ & visual observation \\
$q_t$ & robot proprioception \\
$\ell$ & language instruction \\
$h_t=(o_t,q_t,\ell)$ & policy input context \\
$a_t\in\mathbb{R}^{d_a}$ & robot action \\
$A_{t,H}=(a_t,\ldots,a_{t+H-1})$ & action chunk of horizon $H$ \\
$\theta_k$ & VLA checkpoint after $k$ training steps \\
$\pi_k=\pi_{\theta_k}$ & policy at checkpoint $k$ \\
$\widehat{P}_{\phi}(s_{t+1}\mid s_t,a_t,\ell)$ & learned or simulated world model \\
$\psi_s(s_t)$ & state or observation embedding for comparison \\
\bottomrule
\end{tabular}
\end{center}

The data should be split into training, validation, and test sets:
\begin{equation}
\mathcal{D}=\mathcal{D}_{\mathrm{train}}\cup\mathcal{D}_{\mathrm{val}}\cup\mathcal{D}_{\mathrm{test}}.
\end{equation}
The policy is fine-tuned on $\mathcal{D}_{\mathrm{train}}$, checkpoints are selected using $\mathcal{D}_{\mathrm{val}}$, and final conclusions are reported on $\mathcal{D}_{\mathrm{test}}$.

\section{Action Normalization}

Before computing distances, normalize every action dimension using statistics from the training set:
\begin{equation}
\bar{a}_{t,d} = \frac{a_{t,d}-\mu_d}{\sigma_d+\epsilon},
\qquad d=1,\ldots,d_a,
\end{equation}
where
\begin{equation}
\mu_d = \mathbb{E}_{\mathcal{D}_{\mathrm{train}}}[a_{t,d}],
\qquad
\sigma_d^2 = \operatorname{Var}_{\mathcal{D}_{\mathrm{train}}}(a_{t,d}).
\end{equation}

This is important because translation, rotation, and gripper commands may have very different numerical scales.


\section{\textcolor{red}{What models to be used}}

\begin{enumerate}
    \item  \textbf{\textcolor{red}{World models:} a pretrained light-weighted model: \url{https://rlinf.readthedocs.io/en/latest/rst_source/examples/embodied/opensora.html}}


    \item \textbf{\textcolor{blue}{VLA models:} OpenVLA, Pi0.5-fast, SmoVLA}
\end{enumerate}


\section{Evaluation Mode A: Action-Chunk Evaluation}

For a held-out validation trajectory $\tau_i$ and a window starting at time $t$, define the reference action chunk
\begin{equation}
A^{\star}_{i,t,H}=(a^{\star}_{i,t},a^{\star}_{i,t+1},\ldots,a^{\star}_{i,t+H-1}).
\end{equation}

The model produces a predicted chunk
\begin{equation}
\widehat{A}_{i,t,H}(\theta_k).
\end{equation}

We compare $\widehat{A}_{i,t,H}(\theta_k)$ with $A^{\star}_{i,t,H}$ using DTW, OT, cosine distance, and NLL.

\subsection{One-Step VLA Policy (e.g., OpenVLA)}

For OpenVLA-like models that output only one action at a time, the policy has the form
\begin{equation}
\pi_k(a_t\mid h_t).
\end{equation}

\paragraph{Without a world model}, we form a teacher-forced pseudo-chunk by feeding the held-out state sequence into the model:
\begin{equation}
\widehat{A}^{\mathrm{TF}}_{i,t,H}(\theta_k)
=
(\hat{a}_{i,t},\hat{a}_{i,t+1},\ldots,\hat{a}_{i,t+H-1}),
\end{equation}
where
\begin{equation}
\hat{a}_{i,t+j}=\operatorname{dec}\bigl(\pi_k(\cdot\mid h^{\star}_{i,t+j})\bigr),
\qquad j=0,\ldots,H-1.
\end{equation}
Here $\operatorname{dec}(\cdot)$ denotes the model's action decoder, for example mean action, mode action, or decoded action token. This is a \textbf{teacher-forced action sequence}. It measures action quality on expert states. It does not measure closed-loop compounding error.


\paragraph{With a world model.}
With a world model, a one-step VLA policy can be evaluated in a closed-loop manner.
The policy still outputs only one action at a time. The action chunk is generated by
unrolling the policy together with the learned transition model.

Let the learned world model be
\begin{equation}
\widehat{P}_{\phi}(h_{t+1}\mid h_t,a_t),
\end{equation}
where $h_t$ denotes the policy input at time $t$, including the visual observation,
robot state, language instruction, and any required history.

Starting from the held-out initial state $h^{\star}_{i,t}$, we set
\begin{equation}
\widehat{h}_{i,t}=h^{\star}_{i,t}.
\end{equation}
Then, for $j=0,\ldots,H-1$, we recursively generate
\begin{align}
\widehat{a}_{i,t+j}
&=
\operatorname{dec}\bigl(\pi_k(\cdot\mid \widehat{h}_{i,t+j})\bigr),
\\
\widehat{h}_{i,t+j+1}
&\sim
\widehat{P}_{\phi}(\cdot\mid \widehat{h}_{i,t+j},\widehat{a}_{i,t+j}).
\end{align}

The resulting world-model action chunk is
\begin{equation}
\widehat{A}^{\mathrm{WM}}_{i,t,H}(\theta_k)
=
(\widehat{a}_{i,t},\widehat{a}_{i,t+1},\ldots,\widehat{a}_{i,t+H-1}).
\end{equation}

The corresponding imagined state-action chunk is
\begin{equation}
\widehat{Z}^{\mathrm{WM}}_{i,t,H}(\theta_k)
=
\bigl(
(\widehat{h}_{i,t},\widehat{a}_{i,t}),
(\widehat{h}_{i,t+1},\widehat{a}_{i,t+1}),
\ldots,
(\widehat{h}_{i,t+H-1},\widehat{a}_{i,t+H-1})
\bigr).
\end{equation}

This is different from the teacher-forced pseudo-chunk. In the teacher-forced case,
the model receives the expert states $h^{\star}_{i,t+j}$ at every step. In the
world-model case, only the initial state is taken from the held-out trajectory; all
future states are generated by the world model. Therefore, this setting can measure
closed-loop compounding error.

\subsection{Chunking VLA Policy (Pi 0.5, SmoVLA)}

For chunking models such as $\pi_{0.5}$-style or SmolVLA-style policies, the model directly predicts an action chunk:
\begin{equation}
\pi_k(A_{t,H}\mid h_t).
\end{equation}

The predicted chunk is
\begin{equation}
\widehat{A}^{\mathrm{chunk}}_{i,t,H}(\theta_k)
=
\operatorname{dec}\bigl(\pi_k(\cdot\mid h^{\star}_{i,t})\bigr).
\end{equation}

This is a native chunk prediction from one initial context $h^{\star}_{i,t}$.

\section{Evaluation Mode B: State-Action Rollout with a World Model}

If a world model is available, use it to evaluate the policy in a closed loop.

Start from a held-out initial state:
\begin{equation}
\hat{s}_0=s^{\star}_{i,t}.
\end{equation}

For $j=0,\ldots,H-1$, run
\begin{align}
\hat{h}_j &= \mathcal{O}(\hat{s}_j,\ell_i), \\
\hat{a}_j &= \operatorname{dec}\bigl(\pi_k(\cdot\mid \hat{h}_j)\bigr), \\
\hat{s}_{j+1} &\sim \widehat{P}_{\phi}(\cdot\mid \hat{s}_j,\hat{a}_j,\ell_i).
\end{align}

This produces a generated state-action sequence
\begin{equation}
\widehat{Z}_{i,t,H}(\theta_k)
=
(\hat{z}_0,\hat{z}_1,\ldots,\hat{z}_{H-1}),
\end{equation}
where
\begin{equation}
\hat{z}_j = [\alpha_s\psi_s(\hat{s}_j),\alpha_a\bar{a}_j].
\end{equation}
Here $\alpha_s$ and $\alpha_a$ are fixed weights for state and action features.

The reference state-action chunk is
\begin{equation}
Z^{\star}_{i,t,H}
=
(z^{\star}_{i,t},z^{\star}_{i,t+1},\ldots,z^{\star}_{i,t+H-1}),
\end{equation}
where
\begin{equation}
 z^{\star}_{i,t+j}=[\alpha_s\psi_s(s^{\star}_{i,t+j}),\alpha_a\bar{a}^{\star}_{i,t+j}].
\end{equation}

Then compute DTW, OT, and cosine distance between $\widehat{Z}_{i,t,H}(\theta_k)$ and $Z^{\star}_{i,t,H}$.

\paragraph{Important note.}
The transition model must be written as
\begin{equation}
\widehat{P}_{\phi}(s_{t+1}\mid s_t,a_t,\ell),
\end{equation}
not as $T(s\mid a)$, because the next state depends on the current state, action, and task context.

\section{Metric Definitions}

Let
\begin{equation}
X=(x_0,\ldots,x_{H_x-1}),
\qquad
Y=(y_0,\ldots,y_{H_y-1})
\end{equation}
be two sequences. In action-chunk evaluation, $x_j$ and $y_j$ are normalized actions. In state-action rollout evaluation, they are state-action features.

Use the pointwise cost
\begin{equation}
c(x,y)=\|x-y\|_2^2.
\end{equation}

\subsection{Dynamic Time Warping}

Let $\mathcal{W}$ be the set of valid monotone warping paths between $X$ and $Y$. The DTW distance is
\begin{equation}
D_{\mathrm{DTW}}(X,Y)
=
\frac{1}{|\Gamma^{\star}|}
\min_{\Gamma\in\mathcal{W}}
\sum_{(r,u)\in\Gamma} c(x_r,y_u),
\end{equation}
where $\Gamma^{\star}$ is the minimizing path.

Use DTW when the same behavior may occur at slightly different speeds.

\subsection{Optimal Transport}

Let $p\in\Delta^{H_x}$ and $q\in\Delta^{H_y}$ be uniform distributions over sequence elements. Define
\begin{equation}
\Pi(p,q)=\{\Gamma\in\mathbb{R}_{+}^{H_x\times H_y}:\Gamma\mathbf{1}=p,\Gamma^{\top}\mathbf{1}=q\}.
\end{equation}

The entropic OT distance is
\begin{equation}
D^{\varepsilon}_{\mathrm{OT}}(X,Y)
=
\min_{\Gamma\in\Pi(p,q)}
\sum_{r,u}\Gamma_{ru}c(x_r,y_u)
+
\varepsilon_{\mathrm{ot}}
\sum_{r,u}\Gamma_{ru}(\log\Gamma_{ru}-1).
\end{equation}

If temporal order matters, use the temporally regularized cost
\begin{equation}
c_t(x_r,y_u)
=
\|x_r-y_u\|_2^2
+
\lambda_t\left(\frac{r}{H_x-1}-\frac{u}{H_y-1}\right)^2.
\end{equation}

\subsection{Cosine Distance}

For equal-length sequences, define
\begin{equation}
D_{\mathrm{cos}}(X,Y)
=
1-rac{1}{H}
\sum_{j=0}^{H-1}
\frac{x_j^{\top}y_j}{\|x_j\|_2\|y_j\|_2+\epsilon}.
\end{equation}

Cosine distance measures directional similarity. It should not be used alone because it ignores action magnitude.

\subsection{Log-Likelihood and NLL}

The correct likelihood term is
\begin{equation}
\sum_t \log \pi_k(a_t^{\star}\mid h_t^{\star}),
\end{equation}
not $\sum_t \log(a_t\mid s_t)$.

For a one-step policy under teacher forcing,
\begin{equation}
\mathrm{LL}_{\mathrm{1step}}(k)
=
\sum_{i,t}
\log \pi_k(a^{\star}_{i,t}\mid h^{\star}_{i,t}).
\end{equation}

For a chunking policy with tractable chunk likelihood,
\begin{equation}
\mathrm{LL}_{\mathrm{chunk}}(k)
=
\sum_{i,t}
\log \pi_k(A^{\star}_{i,t,H}\mid h^{\star}_{i,t}).
\end{equation}

Use normalized negative log-likelihood:
\begin{equation}
\mathrm{NLL}(k)
=
-rac{1}{N_{\mathrm{act}}}
\sum \log \pi_k(a^{\star}\mid h^{\star}),
\end{equation}
where $N_{\mathrm{act}}$ is the number of evaluated action tokens or scalar action dimensions.

If the model does not provide a tractable likelihood, report the correct validation surrogate instead:
\begin{itemize}
    \item discrete action-token policy: token NLL;
    \item deterministic continuous policy: L1 or L2 error, or Gaussian-assumed NLL;
    \item diffusion or flow-matching policy: validation diffusion loss or flow-matching loss.
\end{itemize}

\paragraph{World-model rollout likelihood.}
For state-action rollouts, do not use the likelihood of the policy's own sampled actions as the main metric. A bad but overconfident policy can assign high likelihood to its own bad actions.

A better optional metric is matched-reference likelihood. First align the generated rollout with a successful reference rollout using DTW. If $\Gamma^{\star}$ is the alignment path, define
\begin{equation}
\mathrm{LL}^{\mathrm{match}}_{\mathrm{SA}}(k)
=
\sum_{(r,u)\in\Gamma^{\star}}
\log \pi_k(a^{\star}_{u}\mid \hat{h}_{r}).
\end{equation}
This asks whether the policy assigns high probability to the reference action at the generated state.

\section{Aggregating Metrics over Validation Data}

Let $m_{i,t}(k)$ be a lower-is-better metric value for checkpoint $k$ on validation window $(i,t)$. Aggregate over validation windows:
\begin{equation}
\bar{m}(k)
=
\frac{1}{|\mathcal{I}_{\mathrm{val}}|}
\sum_{(i,t)\in\mathcal{I}_{\mathrm{val}}}m_{i,t}(k).
\end{equation}

Use the same validation windows for every checkpoint and every metric.

All metrics should be converted so that lower is better:
\begin{center}
\begin{tabular}{ll}
\toprule
Metric & Lower-is-better form \\
\midrule
DTW & $D_{\mathrm{DTW}}$ \\
OT & $D_{\mathrm{OT}}$ \\
Cosine similarity & $D_{\mathrm{cos}}=1-\mathrm{cos}$ \\
Log-likelihood & $\mathrm{NLL}=-\mathrm{LL}/N_{\mathrm{act}}$ \\
Flow-matching validation objective & $\mathcal{L}_{\mathrm{FM}}$ \\
\bottomrule
\end{tabular}
\end{center}

If combining metrics, normalize each metric first:
\begin{equation}
z_m(k)=
\frac{\bar{m}(k)-\operatorname{median}(m_{\mathrm{ref}})}
{\operatorname{IQR}(m_{\mathrm{ref}})+\epsilon},
\end{equation}
where $m_{\mathrm{ref}}$ is the distribution of the same metric between held-out successful reference demonstrations.

A simple combined score is
\begin{equation}
\mathcal{E}(k)
=
 w_{\mathrm{DTW}}z_{\mathrm{DTW}}(k)
+w_{\mathrm{OT}}z_{\mathrm{OT}}(k)
+w_{\mathrm{cos}}z_{\mathrm{cos}}(k)
+w_{\mathrm{NLL}}z_{\mathrm{NLL}}(k).
\end{equation}
Lower $\mathcal{E}(k)$ is better.

\section{Early-Stopping Rule}

For each metric $m$, select the earliest checkpoint that is close to the best validation value:
\begin{equation}
 k_m^{\star}
 =
 \min\left\{
 k:
 \bar{m}(k)
 \le
 \min_j \bar{m}(j)+\eta_m
 \right\}.
\end{equation}

Here $\eta_m\ge 0$ is a tolerance. A simple choice is
\begin{equation}
\eta_m = 0.01\left|\bar{m}(0)-\min_j \bar{m}(j)\right|.
\end{equation}
This selects the earliest checkpoint within one percent of the best observed metric improvement.

For online training, use patience. Stop when $\bar{m}(k)$ has not improved by at least $\eta_m$ for $P$ consecutive evaluations. Use $P=3$ or $P=5$ as a default.

\section{How to Decide Which Metric Is Best}

To decide which offline metric is best, compare each metric with a trusted downstream quality signal $Q(k)$. Examples of $Q(k)$ are:
\begin{itemize}
    \item simulator success rate;
    \item small real-world validation success rate;
    \item human progress score;
    \item validated world-model success score.
\end{itemize}

For each metric $m$, compute rank correlation:
\begin{equation}
\rho_m = \operatorname{SpearmanCorr}(-\bar{m}(k), Q(k)).
\end{equation}
The negative sign appears because lower metric value should imply higher downstream quality.

Also compute checkpoint-selection regret:
\begin{equation}
\operatorname{Regret}(m)
=
\max_k Q(k)-Q(k_m^{\star}).
\end{equation}

A good early-stopping metric should have high $\rho_m$ and low regret.

If no trusted $Q(k)$ is available, do not claim that one metric is best. You can only report metric smoothness, checkpoint stability, and qualitative agreement.

\section{Experiment Procedure}

\subsection*{Step 1: Prepare Data}

\begin{enumerate}
    \item Split data into training, validation, and test sets.
    \item Compute action normalization statistics using training data only.
    \item Choose the evaluation horizon $H$ and window stride.
    \item Save the same validation windows for all checkpoints.
\end{enumerate}

\subsection*{Step 2: Fine-Tune and Save Checkpoints}

Fine-tune each VLA model and save checkpoints at fixed training intervals:
\begin{equation}
\theta_0,\theta_1,\ldots,\theta_K.
\end{equation}

\subsection*{Step 3: Action-Chunk Evaluation}

For each checkpoint $\theta_k$:
\begin{enumerate}
    \item Generate action chunks on validation windows.
    \item For one-step VLAs, use teacher-forced pseudo-chunks.
    \item For chunking VLAs, use native chunk outputs.
    \item Compute DTW, OT, cosine distance, and NLL or the correct likelihood surrogate.
    \item Aggregate metrics over validation windows.
\end{enumerate}

\subsection*{Step 4: World-Model Rollout Evaluation}

If a world model is available:
\begin{enumerate}
    \item Validate the world model on held-out trajectories.
    \item Roll out each checkpoint from the same held-out initial states.
    \item Construct state-action sequences $\widehat{Z}_{i,t,H}(\theta_k)$.
    \item Compare them with successful reference sequences using DTW, OT, and cosine distance.
    \item Optionally compute matched-reference likelihood.
\end{enumerate}

World-model validation is necessary. At minimum, report the one-step prediction error:
\begin{equation}
E_{\mathrm{WM}}^{1}
=
\mathbb{E}
\left[
\left\|\psi_s(s^{\star}_{t+1})-\psi_s(\hat{s}_{t+1})\right\|_2^2
\right],
\end{equation}
where
\begin{equation}
\hat{s}_{t+1}\sim \widehat{P}_{\phi}(\cdot\mid s^{\star}_t,a^{\star}_t,\ell).
\end{equation}

Do not use rollout horizons longer than the horizon where the world model is reliable.

\subsection*{Step 5: Choose Checkpoints}

For each metric $m$, compute
\begin{equation}
k_m^{\star}.
\end{equation}
Then compare selected checkpoints across metrics.

\subsection*{Step 6: Rank Metrics}

If downstream quality $Q(k)$ is available, rank metrics by:
\begin{enumerate}
    \item Spearman correlation $\rho_m$;
    \item regret $\operatorname{Regret}(m)$;
    \item stability under bootstrap over trajectories or tasks.
\end{enumerate}

\section{Recommended Output Tables}

\subsection*{Metric Curve Table}

\begin{center}
\begin{tabular}{llllr}
\toprule
Model & Checkpoint & Mode & Metric & Value \\
\midrule
OpenVLA& 12200& AC-TF & NLL& 0.71\\
OpenVLA& 13400& AC-TF & OT & 0.73\\
 OpenVLA& 13400& AC-TF & DTW&0.73\\
OpenVLA& 15800& AC-TF & Max iteration& 0.63\\
Chunk VLA & $k$ & AC-Chunk & DTW & \\
Chunk VLA & $k$ & AC-Chunk & OT & \\
Chunk VLA & $k$ & AC-Chunk & NLL/FM & \\
Any VLA & $k$ & SA-WM & DTW & \\
Any VLA & $k$ & SA-WM & OT & \\
Any VLA & $k$ & SA-WM & COS & \\
\bottomrule
\end{tabular}
\end{center}

\subsection*{Metric Ranking Table}

\begin{center}
\begin{tabular}{llllrrr}
\toprule
Metric & Mode & Selected ckpt & External quality & Spearman $\rho$ & Regret & Comment \\
\midrule
DTW & AC-TF & & & & & \\
OT & AC-TF & & & & & \\
COS & AC-TF & & & & & \\
NLL/FM & AC-TF & & & & & \\
DTW & SA-WM & & & & & \\
OT & SA-WM & & & & & \\
COS & SA-WM & & & & & \\
Combined & SA-WM & & & & & \\
\bottomrule
\end{tabular}
\end{center}

\section{Minimal Claims to Make}

The correct claim is:
\begin{quote}
Among the evaluated offline metrics, metric $m$ was the best early-stopping signal because it had the highest rank correlation with downstream quality and the lowest checkpoint-selection regret under our validation protocol.
\end{quote}

Avoid the stronger claim:
\begin{quote}
Metric $m$ proves real-world success without robot evaluation.
\end{quote}

The clean distinction is:
\begin{equation}
\boxed{
\text{teacher-forced action similarity}
\neq
\text{closed-loop state-action quality}
}
\end{equation}

The final early-stopping principle is:
\begin{equation}
\boxed{
\text{choose the earliest checkpoint selected by a validated offline proxy of downstream quality.}
}
\end{equation}

\end{document}
